\section{Lead tracks commercial risk; atypicality does not}\label{sec:limits}

Vocabulary position carries two separable signals, and it is the signed one that tracks
what becomes of a product. Lead marks commercial risk: $+1.4$pp
top-versus-bottom quintile at the use-proof gate with fixed effects,
controls, and owner clustering, and $-0.047$pp per standard deviation in SEC reporting
at equal atypicality, so the filings whose language an industry went on to
adopt belong disproportionately to products that do not survive and to
owners who do not reach public reporting. Atypicality runs the other way and on a different margin: it marks marks that
survive their gate ($+1.3$pp per standard deviation), and its quintile
contrast on SEC reporting is zero to three decimals once class, year and
length are absorbed (Table~\ref{tab:edgar}). A continuous per-standard-deviation
term on the harmonized sample of Table~\ref{tab:staged} does return $+0.042$pp
on a $0.56$\% base; that is the one specification in which it survives, and it
is a linear term through a non-monotone profile.

\dkl{}, adapted from
\citet{Murdock2017, Barron2018} onto topic distributions of USPTO
trademark filings, produces interpretable structure on
commercially-incentivised legal text. The construct is
empirically distinct from patent-track invention within firm: no stable
association survives a change of text representation, and no estimate reaches
a tenth of a standard deviation. And the
corpus supports direct measurement of cross-industry vocabulary
diffusion, with multi-year origin-to-arrival lags, Bass-fittable
adoption, and asymmetric flow out of software on the four classes built so
far.

\dkl{} does not measure innovation in any
strong sense; the defended primitive is vocabulary position on
commercial text, and Appendix~\ref{app:exhibit} shows concretely which
kinds of novelty the primary scoring cannot see. \dkl{} does not
identify an exploitable extreme-tail return premium; the large
debut-specification figure fails a sample-size diagnostic. And the
term-scored signed firm-level advantages of leading filings
(margins, listing) are artifacts of drafting style and lexical
specificity (Appendix~\ref{app:hazard}).

\subsection{Five neighbouring constructs this is not}

\emph{Novelty} and \emph{transience}, in the sense of
\citet{Murdock2017} and \citet{Barron2018}, are the two KL levels
themselves: novelty is divergence from the past, transience divergence
from the future. Atypicality as used here is their average, and so is
neither of them alone; \emph{novelty} in particular is declined because it
names only the backward-looking half of a quantity that faces both ways.
\emph{Resonance} is their signed difference, and is arithmetically what
this paper calls lead. The word is avoided because it names a
mechanism --- the field resonating with an author's contribution --- that
this design cannot observe. A filing can sit where its class was heading
because it was imitated, because it and the class responded to the same
outside event, or by coincidence; the measure does not separate these,
and \emph{resonance} would assert the first.

\emph{Innovation} is a property of the product, not of the sentence
describing it: two firms selling the same thing can draft very differently,
and a genuinely new product can be described in boilerplate
(Appendix~\ref{app:exhibit}). The measure sees prose.

\emph{Radicalness} \citep{DewarDutton1986, TushmanAnderson1986,
HendersonClark1990} is the closest neighbour, since KL is a departure
measure rather than a newness count and the Amazon exhibit is a large
departure built from unremarkable components. It is declined because it
claims that existing competences are devalued, and nothing here licenses the
step from unusual language to devalued competence.

\emph{Distinctiveness} has two prior occupants. In trademark law it is the
registrability doctrine, proven in part by the continued use this paper uses
as an outcome. In strategic management, optimal distinctiveness
\citep{Zhao2017} predicts an inverted-U in conformity --- a shape that
appears here at registration but not at listing, so the term would imply a
theory is being tested when it is not.

What \emph{atypicality} says is narrower than any of them: this description
is unusual for its class and year. Whether the thing described was new,
radical, or good is not measured.

\subsection{Filings are not independent draws}

Some of them share an exact position
with another filing, and outcomes cluster within owner --- the
within-owner correlation is 0.42 for gate survival and is mechanically
1.0 for public listing, since listing is a property of the owner and a
firm with hundreds of marks contributes hundreds of identical successes.
The estimates above are unaffected, because the gate regressions cluster
on normalized owner and the firm-level specifications carry one
observation per owner. But binning filings without that correction
overstates how much structure the corpus contains: on a
positional grid, correcting cell errors for the clustering design effect ---
the factor by which within-owner correlation inflates the variance of a cell
mean --- cuts
apparent excess dispersion by roughly a quarter on the gate outcome and
by half on listing. Displays of the vocabulary plane are descriptive
here for that reason, and the near-zero pooled correlation between $A$ and
$L$ reported in Section~\ref{sec:construct} does not survive cell by cell
within class, so per-cell contrasts should not be read as within-class
effects.

\subsection{What remains unresolved}

What drives the episodic pattern is not settled. Cohort and gate year
move together, so post-crisis bet composition, app-era speculative filings,
and calendar-time enforcement changes cannot be separated on this design.
The representation is also fit on the pooled corpus; a vintage refit would
close the remaining look-ahead channel.

\section{Reproducibility}\label{sec:repro}

\begin{sloppypar}
Code, firm-year panels, and analysis outputs are at
\url{https://github.com/ericsilver/tm-vocabulary}, indexed in the repository
\texttt{README.md}. The pipeline rebuilds from a USPTO Open Data Portal API
key with a single \texttt{make tm} command (the TRTYRAP backfile is about
12 GB, streamed once). Published panels include the SEC-matched firm-year
\dkl{} panel ($n = 59{,}033$) and the PatentsView-joined panel
($n = 174{,}569$). An online appendix under
\texttt{docs/online-appendix} carries a three-panel breakout for each of the
45 Nice classes, the cross-industry comparisons, and a machine-readable
table of per-class contrasts with standard errors.
\end{sloppypar}


